\documentclass[a4paper, serif, 10pt]{moderncv}

%% moderncv
% color and style
\moderncvcolor{black}
\moderncvstyle{banking}

%% page layout/margins
\usepackage[top=2.5cm, bottom=2.5cm, left=1.5cm, right=1.5cm]{geometry}

\nopagenumbers{}

%% Babel config for French language
% ref:
% - https://latex3.github.io/babel/guides/locale-french.html
\usepackage[french]{babel}

%% fontspec
\usepackage{fontspec}

\IfFontExistsTF{Linux Libertine}{
  \setmainfont[Ligatures={TeX, Common}, Numbers=OldStyle]{Linux Libertine}
}{}
\IfFontExistsTF{Linux Libertine O}{
  \setmainfont[Ligatures={TeX, Common}, Numbers=OldStyle]{Linux Libertine O}
}{}

%% CTeX
% CJK support, used to show CJK characters in the resume
%
% - heading=false: disable ctex heading
% - punct=kaiming: use kaiming punctuations styles for CJK
% - scheme=plain: use plain scheme, do not override \normalsize font size
% - space=auto: space settings for CJK characters
%
% CJK fonts are resolved per locale language, e.g. Noto Serif/Sans CJK JP
% for Japanese resumes, with fallback to the next region if not installed
%
% ref:
% - http://ctan.mirrorcatalogs.com/language/chinese/ctex/ctex.pdf
\usepackage[UTF8, heading=false, punct=kaiming, scheme=plain, space=auto]{ctex}

\IfFontExistsTF{Noto Serif CJK SC}{
  \setCJKmainfont{Noto Serif CJK SC}
}{}
\IfFontExistsTF{Noto Sans CJK SC}{
  \setCJKsansfont{Noto Sans CJK SC}
}{}

%% line spacing
\usepackage{setspace}
\setstretch{1.125}

%% URL styling - use normal text instead of monospace
\urlstyle{same}

%% Patch \httplink and \httpslink to handle full URLs
% The original moderncv macros always prepend http:// or https:// to URLs.
% This causes issues when the URL already has a protocol (e.g., https://example.com)
% resulting in malformed URLs like https://https://example.com.
% This patch checks if the URL already contains "://" and uses it directly if so.
% Note: We use \str_set:Nx (x-expansion) to expand macros like \@homepage first.
\ExplSyntaxOn
\str_new:N \g_yamlresume_url_str

\RenewDocumentCommand{\httpslink}{O{}m}{
  \str_set:Nx \g_yamlresume_url_str {#2}
  \str_if_in:NnTF \g_yamlresume_url_str {://}
    { \url{#2} }
    { \url{https://#2} }
}

\RenewDocumentCommand{\httplink}{O{}m}{
  \str_set:Nx \g_yamlresume_url_str {#2}
  \str_if_in:NnTF \g_yamlresume_url_str {://}
    { \url{#2} }
    { \url{http://#2} }
}
\ExplSyntaxOff

\name{Camille Moreau}{}
\title{Ingénieur QA senior spécialisé en automatisation des tests}
\phone[mobile]{+33 6 12 34 56 78}
\email{camille.moreau@example.fr}
\homepage{https://www.camille-moreau-qa.fr}

\address{24 rue des Ateliers, Paris, Île-de-France, France, 75011}{}{}

\extrainfo{{\small \faLinkedin}\ \href{https://www.linkedin.com/in/camille-moreau-qa}{@camille-moreau-qa} {} {} {} • {} {} {} 
{\small \faGithub}\ \href{https://github.com/camillemoreau-qa}{@camillemoreau-qa} {} {} {} • {} {} {} 
{\small \faGitlab}\ \href{https://gitlab.com/camillemoreau-qa}{@camillemoreau-qa}}

\begin{document}

\maketitle

\section{Informations de base}

\cvline{}{Ingénieur QA senior avec plus de 8 ans d'expérience dans la conception, l'automatisation et le pilotage de stratégies de test pour des applications web et des API critiques.
%
\begin{itemize}
\item Expertise en tests fonctionnels, d'intégration, de régression et de performance
\item Maîtrise de Cypress, Playwright, Selenium, Postman, JMeter et des pipelines CI/CD
\item Solide culture qualité : ISTQB, métriques, gestion des risques et amélioration continue
\item Leadership technique, mentorat et collaboration étroite avec les équipes produit et développement
\end{itemize}}
\section{Formation}

\cventry{sept. 2015–juil. 2017}
        {Master, Ingénierie logicielle et qualité, Score : 16/20}
        {Université Paris-Dauphine}
        {\url{https://www.dauphine.psl.eu/}}
        {}
        {\begin{itemize}
\item Formation solide en assurance qualité logicielle et en conduite de projets de test
\item Projets universitaires centrés sur l'automatisation, la non-régression et la mesure de la qualité
\item Pratique des méthodes agiles et de la collaboration avec les équipes de développement
\end{itemize}
\textbf{Cours} : Fondamentaux des tests logiciels, Automatisation des tests fonctionnels, Gestion de la qualité et métriques, Architecture des systèmes d'information, Méthodes agiles et Scrum, Tests de performance et charge}
\section{Expérience professionnelle}

\cventry{mars 2021–Present}
        {Ingénieur QA senior}
        {HexaFinance}
        {\url{https://www.hexafinance.fr}}
        {}
        {\begin{itemize}
\item Pilotage de la stratégie de test multi-produits sur des applications web et des API bancaires critiques
\item Conception et maintenance de suites automatisées avec Cypress, Playwright et REST Assured
\item Réduction de 35 \% des défauts détectés en production grâce à une meilleure couverture des scénarios à risque
\item Mise en place de tests d'API contractuels et de contrôles de régression intégrés aux pipelines GitLab CI
\item Animation de revues de qualité, définition des métriques et accompagnement des équipes produit
\item Mentorat de deux ingénieurs QA juniors et diffusion des bonnes pratiques d'automatisation
\end{itemize}
\textbf{Mots-clés} : Automatisation, Cypress, API REST, CI/CD, Qualité logicielle, Mentorat}

\cventry{janv. 2018–févr. 2021}
        {Ingénieur QA}
        {Novatech Santé}
        {\url{https://www.novatech-sante.fr}}
        {}
        {\begin{itemize}
\item Rédaction des plans de test, des cas de test et des jeux de données pour des applications de santé réglementées
\item Exécution de campagnes de tests fonctionnels, de non-régression et d'intégration
\item Développement de tests automatisés avec Selenium WebDriver et Robot Framework
\item Suivi rigoureux des anomalies dans Jira et coordination des corrections avec les développeurs
\item Participation active aux cérémonies Scrum et aux revues de sprint
\item Contribution à la traçabilité des exigences et à la préparation des livraisons
\end{itemize}
\textbf{Mots-clés} : Tests manuels, Selenium, Jira, Zephyr, SQL, Agile}

\cventry{sept. 2017–déc. 2017}
        {Analyste QA junior}
        {Atelier Logiciel}
        {\url{https://www.atelier-logiciel.fr}}
        {}
        {\begin{itemize}
\item Analyse des exigences fonctionnelles et traduction en scénarios de test vérifiables
\item Création et exécution de tests fonctionnels sur une application de gestion interne
\item Documentation des résultats, des anomalies et des écarts par rapport aux spécifications
\item Collaboration avec les équipes support et développement pour qualifier les incidents
\end{itemize}
\textbf{Mots-clés} : Tests fonctionnels, Postman, Documentation, Reporting}
\section{Langues}

\cvline{Français}{Niveau natif ou bilingue}
\cvline{Anglais}{Niveau professionnel complet \hfill \textbf{Mots-clés} : TOEIC 920, Documentation technique}
\cvline{Espagnol}{Niveau de travail limité \hfill \textbf{Mots-clés} : Conversation professionnelle}
\section{Compétences}

\cvline{Tests automatisés}{Expert \hfill \textbf{Mots-clés} : Cypress, Playwright, Selenium WebDriver, Robot Framework, Cucumber}
\cvline{Tests d'API et intégration}{Avancé \hfill \textbf{Mots-clés} : Postman, REST Assured, Swagger, SoapUI}
\cvline{Tests de performance}{Avancé \hfill \textbf{Mots-clés} : JMeter, Gatling, k6, Analyse de charge}
\cvline{Gestion de la qualité}{Expert \hfill \textbf{Mots-clés} : Stratégie de test, Plans de test, Gestion des anomalies, Métriques qualité, ISTQB}
\cvline{Outils et DevOps}{Avancé \hfill \textbf{Mots-clés} : Jira, Xray, GitLab CI, Docker, Jenkins}
\cvline{Méthodes}{Avancé \hfill \textbf{Mots-clés} : Agile, Scrum, BDD, TDD, Revues de code}
\section{Récompenses}

\cventry{juin 2023}
        {HexaFinance}
        {Prix de l'excellence qualité}
        {}
        {}
        {Récompense attribuée pour la refonte de la stratégie de test et la réduction significative des défauts en production.}
\section{Certificats}

\cventry{avr. 2018}
        {ISTQB}
        {Certification ISTQB niveau fondation}
        {\url{https://www.istqb.org/certifications/foundation-level}}
        {}
        {}

\cventry{nov. 2020}
        {iSQI}
        {Testeur agile certifié}
        {\url{https://www.isqi.org/certified-agile-tester}}
        {}
        {}

\cventry{sept. 2022}
        {ISTQB}
        {ISTQB ingénieur automatisation des tests}
        {\url{https://www.istqb.org/certifications/test-automation-engineer}}
        {}
        {}
\section{Publications}

\cventry{févr. 2023}
        {Automatiser les tests sans sacrifier la maintenabilité}
        {Blog Qualité \& Logiciel}
        {\url{https://www.qualite-logiciel.fr/articles/automatiser-tests-maintenabilite}}
        {}
        {Article pratique sur les bonnes pratiques d'automatisation des tests, la conception de suites robustes et la réduction de la dette de test.}
\section{Projets}

\cventry{janv. 2022–juin 2023}
        {Tableau de bord de pilotage de la qualité logicielle pour les équipes produit.}
        {QualiDash}
        {\url{https://github.com/camillemoreau-qa/qualidash}}
        {}
        {\begin{itemize}
\item Centralisation des métriques de test, des anomalies et des tendances de qualité
\item Visualisation des risques par fonctionnalité et suivi de la couverture automatisée
\item Réduction du temps de préparation des bilans qualité hebdomadaires
\end{itemize}
\textbf{Mots-clés} : Qualité, Métriques, Automatisation, Reporting}

\cventry{mars 2023–déc. 2023}
        {Bibliothèque interne de scénarios réutilisables pour les tests d'API REST.}
        {TestAPIFactory}
        {}
        {}
        {\begin{itemize}
\item Factorisation des appels, assertions et jeux de données pour accélérer l'écriture des tests
\item Intégration dans plusieurs pipelines CI avec exécution parallèle
\item Documentation et formation des équipes à l'utilisation de la bibliothèque
\end{itemize}
\textbf{Mots-clés} : API REST, Tests automatisés, CI/CD, Documentation}
\section{Centres d'intérêt}

\cvline{Technologie}{Automatisation, Logiciel libre, Veille technologique}
\cvline{Sport}{Escalade, Running, Yoga}
\cvline{Culture}{Lecture, Photographie, Cuisine}
\section{Bénévolat}

\cventry{sept. 2020–juin 2023}
        {Bénévole formateur en tests logiciels}
        {Testeurs Solidaires}
        {\url{https://www.testeurs-solidaires.fr}}
        {}
        {\begin{itemize}
\item Animation d'ateliers d'initiation aux tests logiciels pour des personnes en reconversion
\item Accompagnement méthodologique sur la rédaction de cas de test et la gestion des anomalies
\item Création de supports pédagogiques accessibles et adaptés aux débutants
\end{itemize}}
    
\end{document}